% Generated draft: requires mathematical and source review.
\documentclass[11pt]{article}
\usepackage[margin=25mm]{geometry}
\usepackage{fontspec}
\setmainfont{Latin Modern Roman}
\usepackage{amsmath,amssymb,mathtools,unicode-math}
\setmathfont{Latin Modern Math}
\usepackage{enumitem,needspace,xurl,hyperref}
\setlength{\parindent}{0pt}
\setlength{\parskip}{5pt}
\setlength{\emergencystretch}{3em}
\newcommand{\N}{\mathbb{N}}
\newcommand{\Z}{\mathbb{Z}}
\newcommand{\Q}{\mathbb{Q}}
\newcommand{\R}{\mathbb{R}}
\newcommand{\C}{\mathbb{C}}
\newcommand{\F}{\mathbb{F}}
\newcommand{\A}{\mathbb{A}}
\newcommand{\PP}{\mathbb{P}}
\newcommand{\E}{\mathbb{E}}
\newcommand{\eps}{\varepsilon}
\newcommand{\dd}{\,\mathrm{d}}
\begin{document}
\section*{UnsolvedMath ID: 1200045}
\textbf{Source problem code:} catalogue entry
\subsection{Definitions and mathematical statement}
In the standard cost setting, let a countable discrete amenable group \(\Gamma\) act by probability-measure-preserving transformations on a standard probability space \((X,\mu)\). Ergodicity means that every invariant measurable subset has measure \(0\) or \(1\). The action groupoid \(\Gamma\ltimes X\) has arrows \((\gamma,x)\), source \(x\), range \(\gamma x\), and composition \((\delta,\gamma x)(\gamma,x)=(\delta\gamma,x)\). Its groupoid cost is
\[
\operatorname{Cost}(\Gamma\ltimes X)
=\inf_{\{U_i\}}\sum_i\mu(s(U_i)),
\]
where the infimum is over countable measurable partial bisections generating the action groupoid modulo null sets; a partial bisection has injective source and range maps. The action groupoid retains stabilizer arrows, so it must be distinguished from the orbit equivalence relation when the action is not essentially free.

The source asks (i) whether an ergodic essentially faithful amenable action has groupoid cost \(1\), and (ii), if \(\Gamma\) is finitely presented, whether every infinite ergodic action has groupoid cost \(1\). Under the usual faithful-modulo-null-sets interpretation, essential faithfulness means
\[
\{\gamma\in\Gamma:\gamma x=x\text{ for almost every }x\}=\{1\}.
\]
This differs from essential freeness, which requires \(\mu(\{x:\gamma x=x\})=0\) for every \(\gamma\ne1\). The source does not define essential faithfulness or ``infinite action,'' does not explicitly assume probability preservation, and does not explicitly require an infinite group in the first part. In the second part, an intended infinite-orbit interpretation would require \(|\Gamma x|=\infty\) almost everywhere, but that convention cannot be imposed without qualification.

The supplied research summary labels both questions solved and invokes cost \(1\) for ergodic essentially free amenable actions, Ornstein--Weiss, and Gaboriau's fixed-price theory. This does not by itself specify or justify replacing essential faithfulness by essential freeness, treating arbitrary infinite nonfree actions, or omitting the infinitude hypothesis usual in fixed-price-\(1\) statements. Consequently the summary is retained as a supplied assessment; the two literal groupoid-cost targets require source and hypothesis review.
\subsection{Short English statement}
Must the groupoid cost be one for ergodic essentially faithful amenable actions, and for all infinite ergodic actions when the group is finitely presented? The supplied summary discusses essentially free actions, leaving a material gap between its cited theorem and the literal questions.
\subsection{Sources}
\begin{itemize}
\item Local UnsolvedMath registry, permanent ID \(1200045\), exact source code \texttt{catalogue entry}, ``Some Questions --- Question 45.''
\item Abert, \emph{Some questions} (2010), Question 45; \url{https://www.renyi.hu/~abert/questions.pdf}; manual transcription; recorded access 2026-07-29.
\item Supplied research summary names Ornstein--Weiss, Gaboriau's cost, and amenable fixed price \(1\), without complete bibliographic references or a precise theorem covering the nonfree action-groupoid formulations.
\end{itemize}
\end{document}
